\documentclass[12pt]{article}
\usepackage[brazil]{babel}
\usepackage[utf8]{inputenc}
\usepackage[T1]{fontenc}
\usepackage{setspace}
\usepackage{geometry}
\usepackage{xcolor}
\usepackage{tcolorbox}
\geometry{margin=2.5cm}

% Cores
\definecolor{azul}{RGB}{0,102,204}
\definecolor{verde}{RGB}{0,153,76}
\definecolor{roxo}{RGB}{102,0,153}
\definecolor{laranja}{RGB}{255,128,0}

\title{
\textcolor{roxo}{Roteiro de Apresenta\c{c}\~ao Oral}\\
\textbf{\textcolor{azul}{Uso de Nanopart\'iculas e Sensoriamento com Arduino no Monitoramento de Poluentes Aqu\'aticos}}
}

\author{}
\date{}

\begin{document}
\maketitle
\onehalfspacing

\begin{tcolorbox}[colback=blue!5,colframe=azul,title=\textbf{Abertura}]
Bom dia / boa tarde.

Meu nome \'{e} \underline{\hspace{6cm}} e eu vou apresentar um projeto que une \textbf{tecnologia, qu\'imica e meio ambiente}.

O trabalho fala sobre como podemos usar \textbf{Arduino, sensores e nanopart\'iculas} para ajudar a monitorar e melhorar a qualidade da \textbf{\'agua}.
\end{tcolorbox}

\begin{tcolorbox}[colback=green!5,colframe=verde,title=\textbf{Introdu\c{c}\~ao}]
A polui\c{c}\~ao da \textbf{\'agua} \'{e} um problema muito s\'erio, pois afeta rios, lagos e as pessoas que dependem desses recursos.

Pensando nisso, este projeto prop\~oe uma solu\c{c}\~ao simples e acess\'ivel para \textbf{monitorar poluentes aqu\'aticos}, usando sensores eletr\^onicos e nanopart\'iculas.

A ideia principal \'{e} mostrar que a tecnologia pode ajudar na \textbf{sustentabilidade} e tamb\'em no aprendizado dos estudantes.
\end{tcolorbox}

\begin{tcolorbox}[colback=orange!5,colframe=laranja,title=\textbf{Objetivo do Projeto}]
O objetivo geral do projeto \'{e} desenvolver um sistema capaz de \textbf{monitorar e ajudar no tratamento da \'agua polu\'ida}, usando Arduino, sensores e nanopart\'iculas magn\'eticas.

De forma mais espec\'ifica, o projeto busca analisar a qualidade da \textbf{\'agua} e observar como as nanopart\'iculas ajudam a reduzir a polui\c{c}\~ao.
\end{tcolorbox}

\begin{tcolorbox}[colback=purple!5,colframe=roxo,title=\textbf{Sensores Utilizados}]
Para analisar a qualidade da \textbf{\'agua}, foram usados diferentes sensores:

\textbf{Sensor TDS:} mede a quantidade de sais dissolvidos na \'agua, indicando o n\'ivel de polui\c{c}\~ao.

\textbf{Sensor \optico TCRT5000:} detecta mudan\c{c}as na turbidez e na cor da \'agua.

\textbf{Sensor Ultrass\^onico HC-SR04:} auxilia no controle do n\'ivel da \'agua durante os testes.
\end{tcolorbox}

\begin{tcolorbox}[colback=blue!5,colframe=azul,title=\textbf{Como o Sistema Funciona}]
O funcionamento do projeto acontece em tr\^es etapas:

\textbf{1\textordmasculine passo -- Sensoriamento:} os sensores coletam informa\c{c}\~oes da \'agua, como turbidez e condutividade.

\textbf{2\textordmasculine passo -- A\c{c}\~ao das nanopart\'iculas:} as nanopart\'iculas magn\'eticas ajudam a absorver parte dos poluentes.

\textbf{3\textordmasculine passo -- An\'alise:} o Arduino processa os dados e permite acompanhar as mudan\c{c}as na qualidade da \'agua.
\end{tcolorbox}

\begin{tcolorbox}[colback=green!5,colframe=verde,title=\textbf{Resultados}]
Os testes mostraram que o sensor TDS conseguiu identificar claramente o aumento da polui\c{c}\~ao na \'agua.

O sensor \optico tamb\'em respondeu bem \`as mudan\c{c}as de turbidez e cor, principalmente ap\'os a a\c{c}\~ao das nanopart\'iculas.

Esses resultados indicam que o sistema \'{e} eficiente e confi\'avel para o monitoramento da qualidade da \'agua.
\end{tcolorbox}

\begin{tcolorbox}[colback=orange!5,colframe=laranja,title=\textbf{Conclus\~ao}]
Com este projeto, foi poss\'ivel perceber que a uni\~ao entre \textbf{Arduino, sensores e nanopart\'iculas} pode ajudar no cuidado com o meio ambiente.

Al\'em disso, o trabalho contribui para o \textbf{ensino de ci\^encias}, pois envolve conceitos de qu\'imica, f\'isica, programa\c{c}\~ao e sustentabilidade.

O projeto mostra que \textbf{tecnologia acess\'ivel} pode gerar conhecimento e ajudar a enfrentar problemas ambientais reais.

\textbf{Muito obrigado pela aten\c{c}\~ao!}
\end{tcolorbox}

\end{document}

